\chapter{Balanced Diet}

\section*{Definition and Overview}
A balanced diet is one that provides the body with the right kinds and amounts of food to stay healthy. It means eating a wide variety of foods in sensible proportions so the body gets the energy, protein, fats, carbohydrates, vitamins, minerals, and fibre it needs, while avoiding too much added sugar, salt, and unhealthy fat.

In Australia, dietary guidance is built around the five food groups: vegetables and legumes, fruit, grain foods (mostly whole grain), lean meats and poultry, fish, eggs, tofu, nuts and seeds, and milk, yoghurt, cheese, and their alternatives. Foods outside these groups that are high in energy but low in nutrients -- such as soft drinks, lollies, chips, and pastries -- are called discretionary foods and should be limited. A balanced diet supports growth, energy, immunity, and the prevention of chronic disease.

\section*{Epidemiology}
Most Australians do not eat enough vegetables and fruit, and a large share of the energy in the average diet comes from discretionary foods high in sugar, salt, and saturated fat. Around two in three Australian adults are above a healthy weight, and a significant number of children are too. Diet-related conditions -- including heart disease, type 2 diabetes, high blood pressure, dental decay, iron deficiency, and some cancers -- are among the leading causes of illness. At the other end of the spectrum, malnutrition and micronutrient deficiencies also occur, particularly in older people, people with chronic illness, and those with very restricted diets.

\section*{Aetiology and Pathophysiology}
The body needs around 40 essential nutrients, most of which must come from food. Each food group contributes a different mix of nutrients, which is why variety matters:

\begin{itemize}
    \item \textbf{Carbohydrates:} the main energy source; wholegrains and legumes release energy slowly and provide fibre
    \item \textbf{Protein:} needed to build and repair muscle and body tissue; from meat, fish, eggs, dairy, legumes, and nuts
    \item \textbf{Fats:} essential for hormones and cell function, but saturated and trans fats raise cholesterol
    \item \textbf{Vitamins and minerals:} for example iron, calcium, vitamin D, and B vitamins, each with specific roles
    \item \textbf{Fibre and water:} support digestion, gut health, and normal bowel function
    \item \textbf{Energy balance:} weight is gained when energy intake exceeds what the body uses, so portion size matters
\end{itemize}

\section*{Clinical Features}
A poor or unbalanced diet may cause no obvious symptoms for a long time, but signs can include:

\begin{itemize}
    \item Tiredness, poor concentration, and reduced ability to exercise
    \item Unintentional weight loss or slow growth in children
    \item Weight gain, especially around the waist, and difficulty managing weight
    \item Iron deficiency: fatigue, pale skin, breathlessness on exertion
    \item Poor wound healing, frequent infections, or bleeding gums
    \item Tooth decay and gum disease from excess sugar and poor oral care
    \item Constipation from low fibre, or bone and muscle weakness from low calcium and vitamin D
\end{itemize}

\section*{Differential Diagnosis}
\begin{itemize}
    \item Fatigue from an unbalanced diet versus anaemia, thyroid problems, depression, or chronic illness
    \item Weight gain from excess energy intake versus endocrine causes such as an underactive thyroid
    \item Poor growth in children from underfeeding versus malabsorption or underlying medical conditions
    \item Symptoms of deficiency versus medical conditions that cause the same symptoms (for example, breathlessness from anaemia versus heart or lung disease)
    \item Disordered eating versus simple poor dietary habits
\end{itemize}

\section*{When to Seek Medical Care}
A balanced diet is something most people can work on themselves, but seek advice if:

\begin{itemize}
    \item You have unexplained weight loss or rapid weight gain
    \item You have persistent fatigue, weakness, or other symptoms that might be a nutrient deficiency
    \item A child is growing poorly, refusing foods, or falling behind on growth checks
    \item You have a medical condition that requires dietary changes, such as diabetes, coeliac disease, kidney disease, or a food allergy
    \item You are concerned about disordered eating or have struggled to change your eating habits alone
\end{itemize}

A GP can arrange blood tests for deficiencies and refer you to an accredited practising dietitian for personalised advice.

\section*{Diagnosis and Investigations}
\begin{itemize}
    \item \textbf{Diet history:} a food diary or discussion of typical eating patterns, portions, and habits
    \item \textbf{Anthropometry:} weight, height, body mass index (BMI), and waist circumference
    \item \textbf{Growth charts:} in children, plotting weight and height over time
    \item \textbf{Blood tests:} iron studies, vitamin D, vitamin B12, folate, calcium, and thyroid function where deficiency or a medical cause is suspected
    \item \textbf{Risk screening:} blood pressure, blood glucose, and cholesterol checks if weight or diet is a concern
    \item \textbf{Dietitian assessment:} a detailed evaluation of nutrient intake and a tailored plan
\end{itemize}

\section*{Management}
\begin{itemize}
    \item \textbf{Follow the Australian Dietary Guidelines:} eat from the five food groups and limit discretionary foods
    \item \textbf{Vegies and fruit:} aim for a wide variety, with at least two serves of fruit and five of vegetables each day
    \item \textbf{Choose wholegrains:} brown rice, wholemeal bread, and oats instead of refined grains
    \item \textbf{Include lean protein:} lean meat, poultry, fish, eggs, legumes, tofu, and nuts
    \item \textbf{Enjoy dairy or alternatives:} milk, yoghurt, and cheese for calcium and protein
    \item \textbf{Drink water:} make it the main drink and limit sugary drinks, including fruit juice
    \item \textbf{Watch portions:} serve sizes, and be mindful of how much discretionary food is eaten
    \item \textbf{Cook at home:} use a variety of fresh and minimally processed foods, and read nutrition labels on packaged foods
    \item \textbf{Personalised advice:} see an accredited practising dietitian for medical conditions or special diets
\end{itemize}

\section*{Complications}
A poor diet raises the risk of:

\begin{itemize}
    \item Overweight and obesity, and the conditions linked with it, such as type 2 diabetes and joint problems
    \item Heart disease and stroke through excess salt, saturated fat, and sugar
    \item High blood pressure and abnormal cholesterol
    \item Iron deficiency anaemia and other micronutrient deficiencies
    \item Osteoporosis and fractures from low calcium and vitamin D
    \item Dental decay and gum disease
    \item Some cancers and gut problems, such as constipation
    \item In children: poor growth and development; in older adults: frailty and malnutrition
\end{itemize}

\section*{Prognosis and Outcomes}
Improving the balance of the diet can improve energy, weight, and wellbeing, and it lowers the long-term risk of heart disease, diabetes, and other chronic conditions. Changes take time, and sustainable improvements -- rather than strict diets -- are more likely to be kept up. Where deficiencies are present, correcting them with food and, if needed, supplements usually resolves symptoms. In people with medical conditions, dietary change works best as part of overall medical care, and benefits increase the longer healthy habits are maintained.

\section*{Prevention and Self-Care}
\begin{itemize}
    \item Build meals around the five food groups, with plenty of vegetables and fruit
    \item Choose wholegrains and lean proteins most of the time
    \item Limit added sugar, salt, and saturated fat, including in sauces, takeaways, and packaged snacks
    \item Drink water and limit sugary, caffeinated, and alcoholic drinks
    \item Plan and cook meals at home, and read nutrition labels when shopping
    \item Eat regular meals and be mindful of portion sizes
    \item Store and prepare food safely to avoid foodborne illness
    \item See a GP or dietitian if you have health conditions, are concerned about your diet, or want to make significant changes
\end{itemize}

\section*{Sources and Further Reading}
The following sources provide reliable, up-to-date information on healthy eating and a balanced diet.
\begin{itemize}
    \item World Health Organization. \textit{Healthy diet.} https://www.who.int
    \item NHS. \textit{Eating a balanced diet.} https://www.nhs.uk/live-well/
    \item Mayo Clinic. \textit{Nutrition and healthy eating.} https://www.mayoclinic.org
    \item MedlinePlus. \textit{Dietary health.} https://medlineplus.gov
    \item Centers for Disease Control and Prevention. \textit{Nutrition.} https://www.cdc.gov
\end{itemize}
